Set \(h=h_n=n^{-1/4}\), \(\kappa=c_L/4\), \(b_h(r)=(h-r)_+\), and \(a_\Delta(r)=\Delta\phi(r/\Delta)\mathbf1\{0\le r<\Delta\}\), as in the statement.  We first verify that the pair named by \(\mathrm{def:joint\text{-}lower\text{-}handle}\) is legal.  The uniform half-disc design has
\[
 \mathbb P(T=t,R<H)=H^2/2\ge c_0H^2
\]
for \(0<H\le1/2\).  The assignment and consistency restrictions hold by construction.  Also \(|a_\Delta(r)|\le\Delta\), \(0\le b_h(r)\le h\), and \(h+\Delta\le1/2\).  Since \(\kappa\le1/64\),
\[
 |g_{n,\Delta}(r)|\le\kappa(h+\Delta)\le\frac1{128},
\]
so the alternative treated Bernoulli mean lies in \([1/4,3/4]\).

The treated alternative anchor is \(1/2+g_{n,\Delta}(0)\), where
\[
 g_{n,\Delta}(0)=\kappa(h+\Delta).
\]
The map \(b_h\) satisfies \(|b_h(r)-b_h(0)|\le r\).  Since \(\phi'(u)=-3+4u\), one has \(|a_\Delta(r)-a_\Delta(0)|\le3r\) for \(r<\Delta\), while for \(r\ge\Delta\) the same difference equals \(\Delta\le r\).  It follows that
\[
 |g_{n,\Delta}(r)-g_{n,\Delta}(0)|
 \le4\kappa r=c_Lr\le Lr.
\]
All other mean functions are constant.  Hence both laws satisfy \(\mathrm{ass:radial\text{-}lipschitz}\) and belong to \(\mathcal P_L\).

We now use an explicit chi-square calculation to show that the alias adds target separation without adding released divergence.  Let \(Q_0,Q_1\) be the one-unit released laws, define
\[
 \chi^2(Q_1,Q_0)=\int\left(\frac{dQ_1}{dQ_0}-1\right)^2dQ_0,
\]
and let \(J=Q_\Delta(R)/\Delta\).  Absolute continuity holds because every released Bernoulli cell having positive probability under \(Q_1\) has both outcome atoms positive under the baseline \(Q_0\).  Put \(p_j=Q_0(T=1,J=j)\), and, when \(p_j>0\),
\[
 \overline b_j=\mathbb E[b_h(R)\mid T=1,J=j].
\]
The radial density on a treated half-disc, before normalization by its side probability, is \(r\,dr\).  The null identity proved in \(\mathrm{lem:first\text{-}bin\text{-}alias}\) gives
\[
 \int_0^\Delta r a_\Delta(r)\,dr
 =\Delta^3\int_0^1u\phi(u)\,du=0.
\]
The alias vanishes outside the first bin, so under \(Q_1\) the Bernoulli success probability conditional on \((T=1,J=j)\) differs from \(1/2\) by exactly
\[
 \delta_j=\kappa\overline b_j.
\]
Under \(Q_0\), each of the two outcome atoms in that cell has probability \(p_j/2\); under \(Q_1\), their probabilities are \(p_j(1/2+\delta_j)\) and \(p_j(1/2-\delta_j)\).  All control cells agree.  Therefore, with zero-mass cells omitted,
\[
 \chi^2(Q_1,Q_0)
 =4\sum_jp_j\delta_j^2
 \le4\kappa^2\sum_jp_j\mathbb E[b_h(R)^2\mid T=1,J=j]
 =4\kappa^2\int_0^h r(h-r)^2\,dr
 =\frac{\kappa^2h^4}{3},
\]
where the inequality is conditional Jensen and the last integral equals \(h^4/12\).

Assumption \(\mathrm{ass:iid\text{-}sampling}\) makes the released sample laws \(Q_0^{\otimes n}\) and \(Q_1^{\otimes n}\).  Direct multiplication of likelihood ratios gives
\[
 1+\chi^2(Q_1^{\otimes n},Q_0^{\otimes n})
 =(1+\chi^2(Q_1,Q_0))^n
 \le\exp\!\left(\frac{n\kappa^2h^4}{3}\right)
 =\exp(\kappa^2/3).
\]
Cauchy--Schwarz applied to the likelihood ratio gives
\[
 \operatorname{TV}(Q_0^{\otimes n},Q_1^{\otimes n})
 \le\frac12\sqrt{\chi^2(Q_1^{\otimes n},Q_0^{\otimes n})}.
\]
Because \(\kappa\le1/64\) and \(e^x-1\le2x\) for \(0\le x\le1\),
\[
 \operatorname{TV}(Q_0^{\otimes n},Q_1^{\otimes n})
 \le\frac12\sqrt{2/12288}<\frac14.
\]
Since \(\alpha<1/4\), also \(1/4<1-2\alpha\).  Finally the control anchors agree, while the treated anchors differ by \(g_{n,\Delta}(0)\); hence
\[
 \tau_{P_{1,n,\Delta}}(b)-\tau_{P_{0,n,\Delta}}(b)
 =\kappa(h_n+\Delta).
\]
Thus this single same-class pair contains both the sampling displacement \(\kappa h_n\) and the exactly aliased displacement \(\kappa\Delta\), while only the sampling bump enters the divergence.